\section{Related Work}

\subsection{Classical and Deep Survival Models}
Cox proportional hazards~\cite{coxph_original} remains the dominant survival model in clinical
prediction due to its calibration properties and interpretability.
Applied to structured EHR features, CoxPH baselines on large biobank cohorts typically achieve
Harrell C-index in the 0.60--0.68 range for all-cause mortality, driven by biomarkers such as
HbA1c, cystatin~C, and CRP~\cite{harrellcindex}.
Deep survival models, including DeepHit~\cite{deephit} and DeepSurv~\cite{deepsurv}, replace
the linear CoxPH predictor with a neural network but rarely outperform well-regularised CoxPH
on tabular data without imaging or genomic inputs.
Large-scale UK Biobank frameworks have extended this paradigm to joint risk prediction across
thousands of phenotypes~\cite{ukbmdrmf}; respiratory multimorbidity cohorts are identified
as exhibiting the highest structural complexity and the lowest discriminative performance,
motivating the cohort selection and the representation focus of the present study.
All of these tabular and deep tabular approaches operate on pre-extracted fixed-dimension
feature vectors and cannot exploit the sequential and contextual structure present in
longitudinal ICD histories.

\subsection{Temporal Encoding in EHR Models}
Systematic reviews of deep learning on temporal EHR data identify temporal encoding as the
central unsolved methodological challenge~\cite{xie2022temporal}.
Five qualitatively distinct encoding strategies have been proposed, each with a different
notion of what it means to ``represent time.''

\paragraph{Sequence-order encoding.}
The earliest deep EHR models, Doctor~AI~\cite{doctorai} and RETAIN~\cite{retain} (both
NeurIPS 2016), treated the patient history as an ordered sequence of visit vectors processed
by a GRU or reverse-time attention network.
Temporal ordering is preserved, but inter-event duration is invisible: an 8-year gap and
a 1-week gap between the same two diagnoses are indistinguishable if they occupy the same
ordinal position.

\paragraph{Discrete time tokens.}
BERT-style foundation models pretrained on ICD code sequences, including BEHRT~\cite{behrt}
and Med-BERT~\cite{medbert}, encode time through sequence position and bucketed age tokens,
improving over RNNs by modelling long-range co-occurrence but retaining the same ordinal
limitation.
CEHR-BERT~\cite{cehrbert} introduced Artificial Time Tokens (ATTs) inserted between visits
to represent the inter-visit interval in discrete buckets (e.g.\ 1--7 days, 1--4 weeks).
ATTs make duration explicit but remain categorical: two gaps of 8 and 9 years fall in the
same token and are geometrically indistinguishable, a critical limitation for decade-long
mortality trajectories.
CORE-BEHRT and extensions of BEHRT to UK Biobank~\cite{behrtukb} have demonstrated that
this family of models scales well but retains the discrete temporal encoding bottleneck.

\paragraph{Continuous sinusoidal time encoding.}
The sinusoidal positional encoding introduced by Vaswani et al.~\cite{vaswani2017attention}
provides the mathematical foundation for this family of methods, mapping scalar positions
to fixed-frequency sine and cosine components that ensure geometric separation in embedding
space.
Time2Vec~\cite{time2vec} extended this to a learnable, model-agnostic continuous time
representation combining a linear trend component with learnable sinusoidal components,
providing geometric separation between events at different timestamps without discretisation.
Delphi~\cite{delphi2024} applied this design at population scale: a GPT-style
generative transformer trained on $\sim$2 million EHR trajectories uses a fixed-frequency
sinusoidal age encoding to map each event's age-at-occurrence into a continuous geometric
space.
The trajectory embedding in Setting~4 of the present study adopts the same sinusoidal design
principle from Delphi but deploys it within a discriminative CoxPH pipeline, enabling a
controlled comparison against prose serialisation that was not possible in Delphi's generative
framework.

\paragraph{Prose serialisation and LLM encoders.}
The most recent paradigm replaces explicit temporal codes with natural-language text and relies
on the implicit temporal understanding of large pretrained LLMs.
Benchmarking studies have shown that serialising patient records as structured prose and
encoding them with general-purpose LLMs (e.g.\ GTE-Qwen2-7B, LLM2Vec-Llama3.1-8B) frequently
matches or exceeds task-specific EHR foundation models across diverse clinical prediction
tasks~\cite{llmehrenc,serialehr}.
Temporal ordering in these approaches enters through the chronological order of sentences in
the text, but the time gap between events is not geometrically preserved: from surface tokens
alone, the encoder cannot distinguish a 2-week interval from an 8-year one.

\subsection{The Encoding Format Gap}
Despite this progression, no prior study has held the cohort, survival model, and encoder
family constant while systematically varying the temporal encoding format from implicit prose
to continuous geometric representation.
The benchmark established by Harutyunyan et al.~\cite{harutyunyan2019} for ICU time series
and the encoder ablations in the LLM-as-encoder literature~\cite{llmehrenc} vary model
families, not encoding strategies.
The present study isolates encoding format as the sole experimental axis (five serialisation
settings, two encoders, two fusion methods, one CoxPH head), providing the first controlled
evidence that geometric continuity of the time axis, rather than model family or scale, is the
primary structural determinant of long-horizon survival discrimination.
